\section*{अध्याय \ref{ch:b1:organic-redox} --- कार्बनिक रसायन में ऑक्सीकरण और अपचयन}

\begin{solution}{exo:b1:organic-redox:1}
एथेनॉल: \ce{CH3} $-3$, \ce{CH2OH} $-1$। एथेनल: \ce{CH3} $-3$, CHO $+1$।
एथेनोइक अम्ल: \ce{CH3} $-3$, COOH $+3$। प्रोपेनोन: \ce{CH3} $-3$ (दो बार),
C=O $+2$। मेथेनोइक अम्ल: $+2$।
\end{solution}

\begin{solution}{exo:b1:organic-redox:2}
प्रोपेन-1-ऑल: प्रोपेनल, फिर प्रोपेनोइक अम्ल। प्रोपेन-2-ऑल: प्रोपेनोन।
2-मेथिलप्रोपेन-2-ऑल: कोई अभिक्रिया नहीं (नारंगी रंग बना रहता है)।
\end{solution}

\begin{solution}{exo:b1:organic-redox:3}
\ce{CH3COCH3 + 2H+ + 2e- -> CH3CH(OH)CH3};
\ce{CH3CHO + 2H+ + 2e- -> CH3CH2OH}।
\end{solution}

\begin{solution}{exo:b1:organic-redox:4}
\ce{NaBH4}: ब्यूटेन-1-ऑल, ब्यूटेन-2-ऑल, कोई अभिक्रिया नहीं, कोई अभिक्रिया नहीं (अम्ल का
केवल विप्रोटॉनन होता है)। \ce{LiAlH4}: ब्यूटेन-1-ऑल, ब्यूटेन-2-ऑल, ब्यूटेन-1-ऑल और
एथेनॉल, ब्यूटेन-1-ऑल।
\end{solution}

\begin{solution}{exo:b1:organic-redox:5}
\ce{5CH3CH(OH)CH3 + 2MnO4- + 6H+ -> 5CH3COCH3 + 2Mn^2+ + 8H2O}।
$n = 1.20/60.0 = \qty{0.0200}{mol}$; परमैंगनेट $\frac25 \times 0.0200 =
\qty{8.0e-3}{mol}$, अर्थात् \qty{0.020}{mol/L} विलयन का \qty{400}{mL}।
\end{solution}

\begin{solution}{exo:b1:organic-redox:6}
ऐल्कोहॉल को \omterm{def:b1:nernst:couple}{ऑक्सीकारक} के साथ आसवन के लिए सज्जित फ़्लास्क में, दोनों क्वथनांकों के
बीच के ताप पर, गरम कीजिए: प्रोपेनल (\qty{48}{\celsius})
बनते ही आसवित हो जाता है और आगे के ऑक्सीकरण से बच जाता है; प्रोपेन-1-ऑल
(\qty{97}{\celsius}) फ़्लास्क में रहता है।
\end{solution}

\begin{solution}{exo:b1:organic-redox:7}
\ce{BH4-} का एक B--H युग्म कार्बोनिल कार्बन पर आक्रमण करता है, और C=O का $\pi$ युग्म
ऑक्सीजन पर जाता है; \omterm{def:b1:alcohol-activation:alkoxide}{ऐल्कॉक्साइड} एथेनॉल से एक प्रोटॉन ले लेता है। उत्पाद
(\ce{CH3CH2OH}) के कार्बन पर का हाइड्रोजन, C1 पर के दो में से एक, अभिकर्मक से आता है;
ऑक्सीजन पर का हाइड्रोजन विलायक से।
\end{solution}

\begin{solution}{exo:b1:organic-redox:8}
ब्यूटेन-2-ऑल \omterm{def:b1:stereochemistry-in-depth:stereocentre}{किरैल} है, पर समतलीय कीटोन पर दोनों फलकों से समान रूप से आक्रमण
होता है: एक \omterm{def:b1:stereochemistry-in-depth:enantiomers}{रेसेमिक मिश्रण}, प्रकाशतः निष्क्रिय।
\end{solution}

\begin{solution}{exo:b1:organic-redox:9}
\ce{3CH3CH2OH + 2Cr2O7^2- + 16H+ -> 3CH3COOH + 4Cr^3+ + 11H2O}: नारंगी
डाइक्रोमेट (क्रोमियम $+$VI) हरा क्रोमियम(III) बन जाता है।
\end{solution}

\begin{solution}{exo:b1:organic-redox:10}
2-मेथिलब्यूटेन-2-ऑल: एथेनल + \ce{CH3CH2MgBr} ब्यूटेन-2-ऑल देता है (एक
योग, कोई अपचयोपचय नहीं); ब्यूटेनोन तक ऑक्सीकरण; ब्यूटेनोन पर \ce{CH3MgBr},
फिर उपचार। एक ऑक्सीकरण। पेंटेन-3-ऑल: प्रोपेनल + \ce{CH3CH2MgBr},
एक ही चरण में, कोई अपचयोपचय नहीं।
\end{solution}

\begin{solution}{exo:b1:organic-redox:11}
1. कीटोन का चक्रीय \omterm{def:b1:acetals-protection:hemiacetal}{ऐसीटैल} के रूप में \omterm{def:b1:acetals-protection:protecting-group}{रक्षण} कीजिए (एथेन-1,2-डाइऑल, अम्ल,
डीन--स्टार्क)। 2. शुष्क ईथर में \ce{LiAlH4} एस्टर को
प्राथमिक ऐल्कोहॉल (और मेथेनॉल) में अपचित करता है। 3. जलीय अम्ल \omterm{def:b1:acetals-protection:hemiacetal}{ऐसीटैल} को हटाता है:
5-हाइड्रॉक्सीपेंटेन-2-ओन, \ce{CH3CO(CH2)3OH}।
\end{solution}

\begin{solution}{exo:b1:organic-redox:12}
\ce{RCH2OH}: कार्बिनॉल कार्बन $-1$ पर; \ce{RCOOH}: $+3$; चार इलेक्ट्रॉन, जैसे
\ce{RCH2OH + H2O -> RCOOH + 4H+ + 4e-} में। मेथेन ($-4$) से कार्बन डाइऑक्साइड
($+4$): आठ इलेक्ट्रॉन, \ce{CH4 + 2H2O -> CO2 + 8H+ + 8e-}।
\end{solution}

\begin{solution}{pb:b1:organic-redox:1}
\textbf{1.} C1 (OH वहन करने वाला): 0; पाँचों \ce{CH2}: $-2$।
\textbf{2.} C1 (C=O): $+2$; पाँचों \ce{CH2}: $-2$।
\textbf{3.} दोनों COOH कार्बन: $+3$; चारों \ce{CH2}: $-2$।
\textbf{4.} C1 ($0 \to +3$) और उसका पड़ोसी C6 ($-2 \to +3$), जो
दूसरा कार्बोक्सिलिक कार्बन बन जाता है; C1--C6 आबंध टूटता है।
\textbf{5.} $3 + 5 = 8$ इलेक्ट्रॉन।
\textbf{6.} नाइट्रिक अम्ल कार्बोनिल के बगल वाले C--C आबंध को तोड़ने
जितना प्रबल है, जो मृदु \omterm{def:b1:nernst:couple}{ऑक्सीकारक} नहीं करते।
\textbf{7.} \ce{C6H12O + 3H2O -> C6H10O4 + 8H+ + 8e-}।
\textbf{8.} $+$V और $+$I।
\textbf{9.} \ce{2HNO3 + 8H+ + 8e- -> N2O + 5H2O}।
\textbf{10.} जोड़ने पर आठ इलेक्ट्रॉन, \ce{8H+} और जल के तीन अणु
निरस्त हो जाते हैं: \ce{C6H12O + 2HNO3 -> C6H10O4 + N2O + 2H2O}।
\textbf{11.} 8।
\textbf{12.} $10^6/146.0 = \qty{6.85e3}{mol}$ अम्ल; $2 \times 6.85 \times
10^3 \times 63.0 = \qty{8.6e5}{g}$, अर्थात् \qty{0.86}{t} नाइट्रिक अम्ल।
\textbf{13.} \ce{H2O2 + 2H+ + 2e- -> 2H2O}।
\textbf{14.} \ce{C6H12O + 4H2O2 -> C6H10O4 + 5H2O}।
\textbf{15.} $4 \times 6.85 \times 10^3 \times 34.0 = \qty{9.3e5}{g}$,
\qty{0.93}{t}।
\textbf{16.} जल: कोई नाइट्रोजन ऑक्साइड नहीं।
\textbf{17.} बड़ा साम्य स्थिरांक दर के बारे में कुछ नहीं बताता;
\ce{H2O2} \omterm{def:b1:elementary-steps:catalyst}{उत्प्रेरक} के बिना धीरे अभिक्रिया करता है, जैसे उसका अपना अपघटन।
\textbf{18.} वह साइक्लोहेक्सेनोन को वापस साइक्लोहेक्सेनॉल में अपचित करता है:
\ce{4C6H10O + BH4- -> B(OC6H11)4-}, फिर
\ce{B(OC6H11)4- + 4H2O -> 4C6H11OH + B(OH)4-}।
\textbf{19.} \qty{6.85e3}{mol}।
\textbf{20.} \qty{6.85e3}{mol} \ce{N2O} (प्रति अम्ल एक)।
\textbf{21.} $6.85 \times 10^3/0.95 \times 100.0 = \qty{7.2e5}{g}$, \qty{0.72}{t}।
\textbf{22.} $6.85 \times 10^3 \times 44.0 = \qty{3.0e5}{g}$: \textbf{प्रति टन ऐडिपिक अम्ल
\qty{0.30}{t} \ce{N2O}}।
\end{solution}
